% !TeX spellcheck = ru_RU
% !TEX root = vkr.tex

\section*{Введение}
\thispagestyle{withCompileDate}

В современном мире информационные технологии и вычислительные системы стремительно развиваются. Вычислительная мощность устройств увеличивается с каждым годом, а их размеры уменьшаются.

В рамках работы рассматривается задача определения позиций (трекинг) объектов по расстояниям распределенных сенсоров в условиях помех. На фоне технологического прогресса решение такой задачи существенно усложнилось, поскольку объекты наблюдения стали более сложными, быстрыми и непредсказуемыми.

Традиционным подходом к решению задачи является централизованное решение, при котором выделяется единый вычислительный узел, собирающий данные с сенсоров, выполняющий необходимые вычисления и формирующий управляющее воздействие. У такого подхода можно выделить три слабых места:
\begin{itemize}
    \item время на передачу данных от сенсоров до вычислителя;
    \item обработка данных и вычисления;
    \item время на передачу управляющего воздействия от вычислителя обратно к сенсорам.
\end{itemize}
Современные технологии позволяют перейти от этого метода к более новому и эффективному мультиагентному подходу, который устраняет ключевые недостатки централизованной архитектуры, такие как единая точка отказа и ограниченная масштабируемость. В его рамках сенсоры сами выступают в роли вычислительных узлов: они могут самостоятельно определять траекторию объектов, используя как собственные измерения расстояний до целей, так и данные, полученные от соседних сенсоров. Такой подход повышает устойчивость системы, снижает задержки при обработке информации и позволяет более гибко адаптироваться к изменениям в окружающей среде.

К сожалению, на сегодняшний день не существует единой универсальной системы с набором уже реализованных мультиагентных подходов, которая позволяла бы быстро и эффективно протестировать новый мультиагентный алгоритм трекинга и корректно сравнить его с существующими аналогами в одинаковых условиях. Это существенно замедляет проведение исследований, так как исследователям приходится тратить значительное время на самостоятельную реализацию других алгоритмов при отсутствии открытого исходного кода, а также на подготовку симуляций и настройку экспериментальной среды.
